% ---------------------------------------------------------------------------
% trappyscopes -- manuscript draft
%
% TO COMPLETE BEFORE SUBMISSION (kept as comments so they do not print):
%   * Author list and affiliations.
%   * Fleet size and per-instrument configurations.
%   * Totals: recorded experiments, acquisition hours, data volume.
%   * Repository DOI and archived release version; license identifier.
%   * Funding, acknowledgements, data/code availability statements.
%   * A figure pairing the device tree with a rendered experiment record
%     would carry Sections 2-3 far better than prose.
%   * A comparison table against named alternatives would strengthen Sec. 4.
%   * Replace the prose reference note with real numbered citations.
%
% Build:  pdflatex trappyscopes.tex  (twice)
% ---------------------------------------------------------------------------

\documentclass[10pt,a4paper,twocolumn]{article}

\usepackage[T1]{fontenc}
\usepackage[utf8]{inputenc}
\usepackage[margin=1.75cm]{geometry}
\usepackage{listings}
\usepackage{xcolor}
\usepackage[hidelinks]{hyperref}

\definecolor{codebg}{gray}{0.96}
\lstset{
  basicstyle=\ttfamily\tiny,
  backgroundcolor=\color{codebg},
  frame=none,
  xleftmargin=1em,
  framexleftmargin=1em,
  breaklines=true,
  columns=fullflexible,
  keepspaces=true,
  aboveskip=0.5em,
  belowskip=0.5em,
}

\setlength{\parskip}{0.18em}
\setlength{\parindent}{0pt}

\title{\vspace{-1.5em}\textbf{trappyscopes: software-defined laboratory instruments\\with the experiment as a versioned object}\vspace{-0.5em}}

\author{
  Yatharth Bhasin\thanks{Corresponding author: \texttt{yatharth.bhasin@gimm.pt}}\\
  {\small [Co-author names to be completed]}\\
  {\small [Institution and address to be completed]}
}

\date{\small\today}

\linespread{0.97}\selectfont
\begin{document}
\twocolumn[\begin{@twocolumnfalse}
\maketitle

\begin{abstract}
\noindent
We present \texttt{trappyscopes}, a Python framework for building and operating
laboratory instruments as \emph{software-defined assemblies}. An instrument is not
implemented against a fixed driver API; it is declared in a configuration file as a
tree of devices, where each device is an arbitrary Python object --- a camera
binding, a microcontroller running MicroPython, or another networked machine ---
instantiated directly from an import path, with no adapter layer to write. The
framework's principal contribution is that the \emph{experiment}, rather than the
dataset, is its primary versioned object. Every actuation, acquisition, operator
note, protocol checkpoint and file movement is appended to the experiment's own
event chain; each session records the computational environment that produced it;
and the experiment directory can version-control itself, carrying a tracked
manifest of the very files that are too large to commit. Provenance is therefore a
by-product of operating the instrument rather than a curation step performed
afterwards. We describe the architecture, the provenance model, and deployment
across a fleet of parallel microscopes.
\end{abstract}

\vspace{0.4em}
{\small\noindent\textbf{Keywords:} instrument control; research software; provenance;
reproducibility; microscopy; laboratory automation}
\vspace{0.8em}
\end{@twocolumnfalse}]

\section{Introduction}

Two assumptions are built into most instrument-control software, and modern
laboratory practice has outgrown both.

\textbf{The instrument is assumed to be a known object.} Established acquisition
packages expose hardware through a device-adapter layer: a supported device has an
adapter, and an unsupported one is a porting project against the framework's own
API. This is a reasonable trade for standard commercial hardware, but it prices out
the instrument a laboratory actually builds --- a camera with a vendor Python
binding, a pump driven by a microcontroller, an illumination board designed in
house. The integration cost is paid precisely where the novelty is.

\textbf{The computer is assumed to be one computer.} An instrument is increasingly a
small distributed system: a single-board computer at the optics, microcontrollers at
the actuators, a workstation for analysis, networked storage behind them. Frameworks
written around a single host treat everything beyond it as an opaque peripheral
behind a serial link, so a microcontroller's own filesystem, firmware revision and
state fall outside the experimental record entirely.

A third problem compounds both: \textbf{provenance is reconstructed rather than
recorded}. At analysis time the questions are mundane and frequently
unanswerable. Which revision of the acquisition script produced this file? Which
protocol version was followed? Which package versions were installed that afternoon?
Is this directory incomplete, or merely archived? The acquired file carries a
timestamp; the code that made it sits in an uncommitted working tree; the protocol was
a PDF. Electronic notebooks, workflow engines and data-versioning systems each address
part of this, but act \emph{after} acquisition, on whatever the instrument left
behind. What is missing is a record produced \emph{by} running the instrument.

These pressures multiply with parallelism: operating several instruments concurrently
--- the setting this framework was written for --- turns a per-machine configuration
problem into a fleet-management one, in which laboratory-wide policy must coexist with
device-specific settings.

\texttt{trappyscopes} takes the position that an instrument should be \emph{declared}
rather than implemented, and that an experiment should be a self-describing,
self-versioning object rather than a folder of outputs.

\section{Design and implementation}

\subsection{Declarative assembly}
The running instrument is a \texttt{ScopeAssembly}: a tree rooted at the host
machine, whose nodes are devices. The tree is declared, not coded:

\begin{lstlisting}[language=bash]
ScopeAssembly:
  cam:
    kind: detectors.cameras.rpi_hq_picam2.Camera
  pico:
    kind: hive.processorgroups.micropython.SerialMPDevice
    kwargs: {connect: autoconnect, handshake: true}
\end{lstlisting}

\texttt{kind} is a dotted import path that the framework instantiates directly.
There is no adapter layer and no device interface to implement: any importable
Python class is already a device. A vendor binding, a community driver, or a class
written that morning all mount identically, so the instrument is assembled by import
rather than by integration.

Because devices differ in vocabulary, an optional wrapping layer assigns
\emph{roles} without modifying the underlying object: declaring a detector role with
\texttt{read\_method: capture} presents a third-party class through the same
interface as every other detector, with arguments bound at declaration time.
Uniformity is imposed at the edge rather than demanded of the driver.

\subsection{Heterogeneous distribution}
A node in the tree may be another machine. A remote group mounts a second host over
RPyC, after which its devices address as ordinary attributes:

\begin{lstlisting}[language=Python]
scope.MS1.read()       # sensor on this machine
scope.M2.MS2.read()    # sensor on another machine
\end{lstlisting}

MicroPython boards are first-class members of the same tree rather than serial
peripherals behind it. The framework manages a board's own filesystem --- mirroring
firmware, skipping unchanged files, honouring a manifest declared by the firmware
repository --- so microcontroller code falls under the same version control and record
as everything else.

\subsection{Fleet configuration}
A device's configuration may be extended by laboratory-wide files. Most keys
\emph{replace}: the highest-priority source wins outright, correct for per-device
settings such as a data directory. A small, documented set instead \emph{expands},
unioning values across layers, so a laboratory manifest can add a shared protocol
repository or a recognised operator without overwriting what an instrument declares for
itself. Which keys expand is a statement about what each key means, recorded per key
rather than inferred from its shape.

\section{The experiment as a versioned object}

An experiment is a directory, qualified by a marker file rather than a database entry.
It can be read, copied, archived or inspected with ordinary tools, and carries its own
record.

\subsection{A single event chain}
Every occurrence is appended to one chronological chain. Acquisitions, operator
notes, protocol checkpoints, scheduled-task registrations, script executions and
data transfers are the same kind of record, each stamped with instrument identity,
session, experiment identifier, wall-clock time, elapsed experiment time and a
monotonic machine time:

\begin{lstlisting}[language=bash]
- type: cam_acq_img
  filename: chamber10.png
  dt: 2025-09-29 16:17:32.973482
  exptime: 652.70
  eid: 9b400bd83c
  scopeid: VWR
- type: user_note
  note: Focus set at Chamber 9. Moving towards 2.
\end{lstlisting}

Because an operator note and a machine acquisition are the same type of object, the
narrative and the data share one timeline. Measurement streams build on the same record:
typed, self-describing series accumulating into a dataframe while carrying per-field
units, labels and descriptions.

\subsection{Session-scoped environment capture}
A \emph{session} is one sitting at the instrument, and its role is to capture what
produced the data. At session start the framework records the operator, the commit
of the control software itself, the complete list of installed Python distributions
with versions, the interpreter version, and the host platform and identity. These are written into the experiment, so ``what code and what environment produced
this?'' is answered by the experiment rather than from memory. An unchanged environment
is stored once and referenced thereafter; a change --- an upgraded package, a different
machine, a new commit --- produces a new full record, precisely the event worth
noticing.

\subsection{Protocol identity}
Protocols are resolved from a version-controlled repository. On loading, the
framework records the commit, whether the working tree was clean, and a permalink to
that exact revision. The protocol that ran is therefore \emph{identified} rather than
named --- the difference between ``we followed the staining protocol'' and ``we
followed this revision of it, and it carried no uncommitted edits.''

\subsection{Self-versioning directories}
An experiment directory may optionally be its own version-control repository,
committing when opened and closed, tagged with the session, so its metadata's evolution
is inspectable with ordinary tooling.

Bulk acquisition data is deliberately excluded from that repository: microscopy video
and image stacks do not belong in commit history. This creates an obvious gap, which
the framework closes with a tracked manifest of the directory that \emph{includes}
the excluded files. The history of one small text file is therefore a readable record
of what existed and how the experiment's structure changed, even where the content
itself was never committed.

\subsection{Two legible copies}
Acquisition machines fill up, so data must move. The framework classifies each file by
what it is: metadata --- identity marker, event log, session record, protocol payload,
scripts --- is \emph{copied}, while bulk data is \emph{moved}. Both sides remain
complete, legible experiments: storage holds everything, and the acquisition machine
retains a full record of what happened with the bulk data removed.

The depleted copy then declares itself. A \emph{policy} is a durable statement about
an experiment, emitted both as an event and as a marker file whose \emph{name} carries
the signal, so a tool recognises it from a directory listing with no parsing at all. A
\texttt{stub} policy answers ``is this incomplete, or simply archived?'' immediately.
The mechanism generalises: a \texttt{contaminated} policy marks an experiment whose
course or data structure changed materially, so downstream analysis meets that fact
before reading anything. Policies are scoped: \texttt{stub} describes one copy and
never travels, while \texttt{contaminated} is true of the experiment everywhere and
always does.

\section{Novelty and relation to existing work}

The individual mechanisms are not each unprecedented; the contribution is the position
they are taken from and the fact that they compose.

\textbf{Provenance as a by-product of operation.} This is the central claim. Code
identity, computational environment, protocol revision, operator, instrument
configuration and data location are captured because the instrument ran, not because
someone later assembled them. Workflow engines, data-versioning tools and electronic
notebooks address parts of this problem, but operate on the output of acquisition.
\texttt{trappyscopes} places the record at the instrument, during acquisition, in the
same object the data lands in.

\textbf{The experiment directory as a self-versioning object.} We are not aware of
another instrument-control framework in which the experiment directory is itself a
repository that commits on open and close and carries a tracked manifest covering the
files version control excludes. That manifest gives structural history to data too
large to version, and is to our knowledge specific to this work.

\textbf{A zero-adapter device model.} Declaring a device as an import path, and
optionally assigning a role through method binding, removes the adapter-writing step
entirely. The significance is coverage rather than convenience: self-built instruments
become first-class without a porting effort, which is where adapter-based
architectures impose their highest cost.

\textbf{A heterogeneous device tree.} Single-board computers, workstations and
microcontrollers are addressed through one uniform tree, with the microcontroller's
filesystem managed and versioned rather than opaque. Distributed instrument control
exists; folding the microcontroller into the same naming and provenance regime is less
common. Finally, the two-copy model with a self-declared stub gives a machine-readable
answer to a question normally resolved by asking a colleague.

\subsection*{What is not claimed}
No new imaging modality, optical design or analysis method is proposed, and no claim is
made of faster acquisition or better image quality. We do not claim general superiority
over mature acquisition software: those systems target broad commercial-hardware support
and sophisticated acquisition control, whereas this framework targets heterogeneous,
self-built, parallel instruments and provenance-first operation; for standard commercial
hardware they remain the better choice. The provenance model records what the framework
mediates \textemdash{} physical actions outside it are captured only if an operator notes
them.

\section{Availability and deployment}

\texttt{trappyscopes} is implemented in pure Python (3.12 or later) and runs on
Raspberry Pi OS, macOS and Linux. A device is configured by a single file, optionally
layered with laboratory-wide ones. The operator-facing surface is an interactive Python
session in which the assembly, the open experiment and the user tools are bound as
ordinary names, so acquisition is scripted against live objects rather than through a
separate macro language. Scripts executed through the framework are copied into the
experiment as they run, versioned when they change, and recorded as events, so the code
that produced a result travels with it.

Assembly and device tree, the event and session model, measurement streams, protocol
resolution, scripting with payload capture, per-experiment version control and
synchronisation to laboratory storage are in routine use. Recorded experiments in the
reference deployment date from 2025 onward, on Raspberry Pi--based microscopes operating
in parallel alongside a development workstation. Live plotting and a restructuring of
the assembly layer remain under development.

\section{Limitations}

The framework tracks a single current experiment and binds the process working
directory to it, so parallelism is across machines rather than within a session. File
classification is per top-level entry, so data nested inside an otherwise-metadata
directory is copied with it rather than moved. Synchronisation currently shells out to
an external binary outside the Python dependency set, which must be installed per
machine; a pure-Python client would keep the dependency story within the package
manager at the cost of reimplementing transfer logic. We report no acquisition-latency
or overhead measurements: the provenance machinery writes small per-event records and
is not expected to be significant against image acquisition, but this is unmeasured and
is not asserted. Windows is untested.

\section{Conclusion}

Treating the experiment rather than the dataset as the primary versioned object changes
what a laboratory has to remember. Code identity, environment, protocol revision and
data location become properties the experiment carries, recorded as a consequence of
running the instrument. Combined with a declarative, adapter-free device model that
admits self-built and distributed hardware on equal terms, this offers a practical
route to reproducible operation for the instruments laboratories increasingly build for
themselves.

\section*{References}
\footnotesize
To be completed with full bibliographic detail. The argument should engage:
Micro-Manager and Pycro-Manager; workflow systems (Snakemake, Nextflow); data
versioning (DVC, git-annex); provenance standards (W3C PROV, RO-Crate); the FAIR
principles; the electronic laboratory notebook literature; and open-hardware
microscopy (e.g.\ OpenFlexure).

\end{document}
